%% title: Heteroskedasticity (WLS) and AR(1) Errors (GLS)
\documentclass[11pt]{article}
\usepackage{metricsstyle}

\begin{document}

\lessonheader{8}{Weighted least squares, AR(1) errors, and GLS with a known $\rho$}{}

The setting is the one Lesson 7 ended with:
\[ y = X\beta + u, \qquad u \sim N(0, \sigma^{2}\Omega), \qquad
   \Omega = \begin{bmatrix}
     w_{11} & w_{12} & \cdots & w_{1n}\\
     w_{21} & w_{22} & \cdots & w_{2n}\\
     \vdots & \vdots & \ddots & \vdots\\
     w_{n1} & w_{n2} & \cdots & w_{nn}
   \end{bmatrix}. \]
GLS needs a matrix $T$ with $T'T = \Omega^{-1}$ (or a scalar multiple of it); then
OLS on $y_{*} = Ty$, $X_{*} = TX$ gives $b_{*} = (X'\Omega^{-1}X)^{-1}X'\Omega^{-1}y$.
This lesson finds $T$ in two special cases.


\section{Case 1: heteroskedasticity and weighted least squares}

\topic{The variance depends on a variable $z_{i}$}
The errors are uncorrelated, but their variance moves with $z_{i}$:
\[ \Var(u_{i}) = \sigma^{2}w_{ii}, \qquad w_{ii} = h(z_{i}) > 0, \qquad
   \Omega = \begin{bmatrix}
     h(z_{1}) & 0 & \cdots & 0\\
     0 & h(z_{2}) & \cdots & 0\\
     \vdots & \vdots & \ddots & \vdots\\
     0 & 0 & \cdots & h(z_{n})
   \end{bmatrix}, \]
\[ y_{i} = \beta_{1} + \beta_{2}X_{2i} + \beta_{3}X_{3i} + \cdots + \beta_{K}X_{Ki} + u_{i},
   \qquad u_{i} \sim N\bigl(0, \sigma^{2}h(z_{i})\bigr). \]

\topic{Example: $h(z_{i}) = z_{i}^{2}$}
\[ u_{i} \sim N(0, \sigma^{2}z_{i}^{2})
   \quad\Longrightarrow\quad
   u_{*i} = \frac{u_{i}}{z_{i}} \sim N(0, \sigma^{2}). \]
In general,
\[ u_{*i} = \frac{u_{i}}{\sqrt{h(z_{i})}} \sim N(0, \sigma^{2}). \]

\topic{Weighted least squares}
Divide every term of the model by $\sqrt{h(z_{i})}$:
\[ \underbrace{\frac{y_{i}}{\sqrt{h(z_{i})}}}_{y_{*i}}
   = \beta_{1}\frac{1}{\sqrt{h(z_{i})}}
     + \beta_{2}\frac{X_{2i}}{\sqrt{h(z_{i})}}
     + \cdots
     + \beta_{K}\frac{X_{Ki}}{\sqrt{h(z_{i})}}
     + \underbrace{\frac{u_{i}}{\sqrt{h(z_{i})}}}_{u_{*i}}. \]
The transformed errors are homoskedastic, so OLS on the transformed data is GLS.
Note that the constant becomes the regressor $1/\sqrt{h(z_{i})}$, so the transformed
regression has no intercept. In the notation of Lesson 7, each observation is
weighted by the inverse of its standard deviation:
\[ T = \begin{bmatrix}
     1/\sqrt{h(z_{1})} & 0 & \cdots & 0\\
     0 & 1/\sqrt{h(z_{2})} & \cdots & 0\\
     \vdots & \vdots & \ddots & \vdots\\
     0 & 0 & \cdots & 1/\sqrt{h(z_{n})}
   \end{bmatrix}, \qquad T'T = \Omega^{-1}. \]


\section{Case 2: autocorrelation}

The second special case of Lesson 7 has a common variance $\bar w$ on the diagonal and
non-zero covariances off it:
\[ \Omega = \begin{bmatrix}
     \bar w & w_{12} & \cdots & w_{1n}\\
     w_{21} & \bar w & \cdots & w_{2n}\\
     \vdots & \vdots & \ddots & \vdots\\
     w_{n1} & w_{n2} & \cdots & \bar w
   \end{bmatrix}. \]

\topic{The AR(1) model}
With time series data, indexed by $t = 1, \dots, n$,
\[ y_{t} = \underbrace{x_{t}'}_{1\times K}\,\underbrace{\beta}_{K\times 1} + u_{t},
   \qquad u_{t} = \rho u_{t-1} + \varepsilon_{t}, \qquad |\rho| < 1, \]
where the innovations $\varepsilon_{t}$ are white noise:
\[ \E(\varepsilon_{t}) = 0 \ \ \forall t, \qquad
   \E(\varepsilon_{t}\varepsilon_{s}) = 0 \ \ (t \neq s), \qquad
   \varepsilon \sim N(0, \sigma_{\varepsilon}^{2}I_{n}). \]

\topic{The lag operator}
\[ Lx_{t} = x_{t-1}, \qquad L^{2}x_{t} = L(Lx_{t}) = x_{t-2}, \qquad
   L^{p}x_{t} = x_{t-p}. \]
The operator acts on what stands to its right, so
\[ u_{t} = \rho Lu_{t} + \varepsilon_{t}
   \quad\Longrightarrow\quad (1 - \rho L)\,u_{t} = \varepsilon_{t}
   \quad\Longrightarrow\quad u_{t} = (1 - \rho L)^{-1}\varepsilon_{t}. \]

\topic{Math recall: the geometric series}
\begin{align*}
  S &= 1 + a + a^{2} + \cdots\\
  aS &= \phantom{1 + {}} a + a^{2} + \cdots\\
  S - aS &= 1 \quad\Longrightarrow\quad S = \frac{1}{1 - a} \qquad (|a| < 1).
\end{align*}

\topic{Connect to $u_{t}$: a moving average of infinite order}
With $a = \rho L$ and $|\rho| < 1$,
\begin{align*}
  u_{t} &= \bigl(1 + \rho L + \rho^{2}L^{2} + \cdots\bigr)\varepsilon_{t}\\
        &= \varepsilon_{t} + \rho\varepsilon_{t-1} + \rho^{2}\varepsilon_{t-2} + \cdots
         = \sum_{j=0}^{\infty}\rho^{j}\varepsilon_{t-j}
         \qquad \text{MA}(\infty).
\end{align*}
In particular $u_{t-1}$ depends only on $\varepsilon_{t-1}, \varepsilon_{t-2}, \dots$,
so it is uncorrelated with $\varepsilon_{t}$:
\[ \E(u_{t-s}\,\varepsilon_{t}) = 0 \qquad \text{for all } s \geq 1. \]

\topic{We need to find $\Omega$: the variance}
Square $u_{t} = \rho u_{t-1} + \varepsilon_{t}$ and take expectations:
\begin{align*}
  u_{t}^{2} &= \rho^{2}u_{t-1}^{2} + \varepsilon_{t}^{2} + 2\rho\,u_{t-1}\varepsilon_{t},\\
  \E(u_{t}^{2}) &= \rho^{2}\E(u_{t-1}^{2}) + \E(\varepsilon_{t}^{2})
                   + 2\rho\underbrace{\E(u_{t-1}\varepsilon_{t})}_{=\,0}.
\end{align*}
The variance is the same at every $t$, $\E(u_{t}^{2}) = \E(u_{t-1}^{2}) = \sigma_{u}^{2}$,
so
\[ \sigma_{u}^{2} = \rho^{2}\sigma_{u}^{2} + \sigma_{\varepsilon}^{2}
   \quad\Longrightarrow\quad \sigma_{u}^{2}(1 - \rho^{2}) = \sigma_{\varepsilon}^{2}
   \quad\Longrightarrow\quad \sigma_{u}^{2} = \frac{\sigma_{\varepsilon}^{2}}{1 - \rho^{2}}. \]

\topic{The covariances}
Multiply $u_{t} = \rho u_{t-1} + \varepsilon_{t}$ by $u_{t-s}$ and take expectations;
the $\varepsilon_{t}u_{t-s}$ term always drops out.
\begin{align*}
  \text{one lag:}\quad
  \E(u_{t}u_{t-1}) &= \rho\E(u_{t-1}^{2}) + \E(\varepsilon_{t}u_{t-1})
                    = \rho\sigma_{u}^{2},\\
  \text{two lags:}\quad
  \E(u_{t}u_{t-2}) &= \rho\underbrace{\E(u_{t-1}u_{t-2})}_{\rho\sigma_{u}^{2}}
                      + \E(\varepsilon_{t}u_{t-2})
                    = \rho^{2}\sigma_{u}^{2},\\
  \text{three lags:}\quad
  \E(u_{t}u_{t-3}) &= \rho\underbrace{\E(u_{t-1}u_{t-3})}_{\rho^{2}\sigma_{u}^{2}}
                      + \E(\varepsilon_{t}u_{t-3})
                    = \rho^{3}\sigma_{u}^{2}.
\end{align*}
In general
\[ \E(u_{t}u_{t-s}) = \rho^{s}\sigma_{u}^{2}, \qquad s = 0, 1, 2, \dots \]

\topic{The AR(1) covariance matrix}
Entry $(t,s)$ is $\sigma_{u}^{2}\rho^{|t-s|}$: the correlation dies out geometrically
away from the diagonal.
\[ \underbrace{\Var(u)}_{n\times n} = \sigma_{u}^{2}\,\Omega, \qquad
   \Omega = \begin{bmatrix}
     1          & \rho       & \rho^{2}   & \cdots & \rho^{n-1}\\
     \rho       & 1          & \rho       & \cdots & \rho^{n-2}\\
     \rho^{2}   & \rho       & 1          & \cdots & \rho^{n-3}\\
     \vdots     & \vdots     & \vdots     & \ddots & \vdots\\
     \rho^{n-1} & \rho^{n-2} & \rho^{n-3} & \cdots & 1
   \end{bmatrix}. \]
So $\sigma^{2} = \sigma_{u}^{2}$ and $\bar w = 1$ in the opening notation, and $\Omega$
depends on the single unknown $\rho$. GLS needs $T$ with $T'T \propto \Omega^{-1}$;
the next section builds it.


\section{How do we get GLS?}

\topic{(1) $\rho$ is known}
Write the model at $t$, and the model at $t-1$ multiplied by $\rho$:
\begin{align*}
  y_{t}        &= x_{t}'\beta + u_{t}, \tag{1}\\
  \rho y_{t-1} &= \rho x_{t-1}'\beta + \rho u_{t-1}. \tag{2}
\end{align*}
Then (1) $-$ (2):
\[ y_{t} - \rho y_{t-1} = (x_{t} - \rho x_{t-1})'\beta
   + \underbrace{u_{t} - \rho u_{t-1}}_{\varepsilon_{t}\ \text{iid}},
   \qquad t = 2, \dots, n, \]
that is, with the lag operator,
\[ (1 - \rho L)\,y_{t} = \bigl[(1 - \rho L)\,x_{t}\bigr]'\beta + \varepsilon_{t}. \]
Call the quasi-differenced data $y_{*t} = y_{t} - \rho y_{t-1}$ and
$x_{*t} = x_{t} - \rho x_{t-1}$; then $y_{*t} = x_{*t}'\beta + \varepsilon_{t}$ has
iid errors, and OLS on it is GLS.

\topic{The quasi-differencing matrix $T_{*}$}
In matrix form, dropping the first observation (there is no $y_{0}$):
\[ \underbrace{T_{*}}_{(n-1)\times n} = \begin{bmatrix}
     -\rho  & 1      & 0      & \cdots & 0      & 0\\
     0      & -\rho  & 1      & \cdots & 0      & 0\\
     \vdots & \vdots & \ddots & \ddots & \vdots & \vdots\\
     0      & 0      & \cdots & -\rho  & 1      & 0\\
     0      & 0      & \cdots & 0      & -\rho  & 1
   \end{bmatrix}, \]
\[ y_{*} = T_{*}y = T_{*}\begin{bmatrix} y_{1}\\ y_{2}\\ \vdots\\ y_{n} \end{bmatrix}
   = \begin{bmatrix} y_{2} - \rho y_{1}\\ y_{3} - \rho y_{2}\\ \vdots\\ y_{n} - \rho y_{n-1} \end{bmatrix},
   \qquad X_{*} = T_{*}X. \]
$T_{*}'T_{*}$ is close to, but not exactly, proportional to $\Omega^{-1}$ (see below).

\topic{The exact $T$: keeping the first observation}
\[ \underbrace{T}_{n\times n} = \begin{bmatrix}
     \sqrt{1 - \rho^{2}} & 0      & 0      & \cdots & 0\\
     -\rho               & 1      & 0      & \cdots & 0\\
     0                   & -\rho  & 1      & \cdots & 0\\
     \vdots              & \vdots & \ddots & \ddots & \vdots\\
     0                   & 0      & \cdots & -\rho  & 1
   \end{bmatrix}
   \quad \longleftarrow \text{the 1st row handles the 1st observation,} \]
\[ Ty = \begin{bmatrix} \sqrt{1 - \rho^{2}}\,y_{1}\\ y_{2} - \rho y_{1}\\ \vdots\\ y_{n} - \rho y_{n-1} \end{bmatrix}. \]
The first transformed error is $\sqrt{1 - \rho^{2}}\,u_{1}$, with variance
$(1 - \rho^{2})\sigma_{u}^{2} = \sigma_{\varepsilon}^{2}$, the same as every
$\varepsilon_{t}$. So $\Var(Tu) = \sigma_{\varepsilon}^{2}I_{n}$, and
\[ T'T = (1 - \rho^{2})\,\Omega^{-1} = \begin{bmatrix}
     1      & -\rho      & 0          & \cdots & 0      & 0\\
     -\rho  & 1+\rho^{2} & -\rho      & \cdots & 0      & 0\\
     0      & -\rho      & 1+\rho^{2} & \cdots & 0      & 0\\
     \vdots & \vdots     & \vdots     & \ddots & \vdots & \vdots\\
     0      & 0          & 0          & \cdots & 1+\rho^{2} & -\rho\\
     0      & 0          & 0          & \cdots & -\rho  & 1
   \end{bmatrix}, \]
which is $\Omega^{-1}$ up to a scalar; the scalar does not change
$b_{*} = (X'\Omega^{-1}X)^{-1}X'\Omega^{-1}y$. The matrix $T_{*}'T_{*}$ is the same
except that its $(1,1)$ entry is $\rho^{2}$ instead of $1$: dropping the first
observation loses exactly that piece.


\section{The AR(2) process}

\[ u_{t} = \rho_{1}u_{t-1} + \rho_{2}u_{t-2} + \varepsilon_{t}, \qquad
   \rho_{1}, \rho_{2} \text{ known}. \]
With the lag operator,
\begin{align*}
  & (1 - \rho_{1}L - \rho_{2}L^{2})\,u_{t} = \varepsilon_{t}\\
  \Longleftrightarrow\quad & u_{t} - \rho_{1}Lu_{t} - \rho_{2}L^{2}u_{t} = \varepsilon_{t}\\
  \Longleftrightarrow\quad & u_{t} - \rho_{1}u_{t-1} - \rho_{2}u_{t-2} = \varepsilon_{t}.
\end{align*}


\begin{transcriptnotes}
  \item The set opens with a divider page carrying only ``Lesson 8''; it is not
        reproduced. The lecture stops after the AR(2) definition on page 7.
  \item As in Lessons 3--7, $6^{2}$ is $\sigma^{2}$ (and $6_{u}^{2}$,
        $6_{\varepsilon}^{2}$ are $\sigma_{u}^{2}$, $\sigma_{\varepsilon}^{2}$); $B$
        for $\beta$ is typeset $\beta$; $k$ for the number of regressors is typeset
        $K$. The regressor row, written $X_{t}'$, is typeset $x_{t}'$ (a $1\times K$
        row of $X$). $T^{*}$ is typeset $T_{*}$ to match $b_{*}$, $X_{*}$ of
        Lesson 7.
  \item Notation (layout revision): transformed variables are written
        $y_{*i}$, $u_{*i}$, $y_{*t}$, $x_{*t}$ (the page has $u_{i}^{*}$) to match
        $y_{*} = Ty$ of Lesson 7. The lag operator is written in front of what it
        acts on, $(1-\rho L)u_{t}$ and $(1-\rho L)^{-1}\varepsilon_{t}$, where the page
        has $u_{t}(1-\rho L)$ and $\varepsilon_{t}/(1-\rho L)$. \textbf{Corrected:}
        the page writes $w_{ii} = \sigma^{2}h(z_{i})$, which counts $\sigma^{2}$ twice
        given $\Var(u) = \sigma^{2}\Omega$; typeset $w_{ii} = h(z_{i})$. The regressor
        row $x_{t}$ is lower case so that $Lx_{t} = x_{t-1}$ does not clash with the
        matrix $X$.
  \item Layout (this revision): every matrix is written out with its full corner
        pattern and stands in its own display; the sentence on why GLS needs $T$,
        the diagonal $\Omega$ and $T$ of Case 1, the remark on the intercept, the
        explicit $Ty$, the matrix $(1-\rho^{2})\Omega^{-1}$, and the $(1,1)$-entry
        comparison with $T_{*}'T_{*}$ are \textbf{added}. The covariance derivations
        at lags 1--3 are merged into one aligned block; the products
        $u_{t}u_{t-s} = \ldots$ that precede each expectation on the page are folded
        into the sentence above it.
  \item Page 2: the opening $\Omega$ and the line $u \sim N(0, \sigma^{2}\Omega)$ are
        the recap from Lesson 7; the sentence introducing them is \textbf{added}.
        \textbf{Corrected:} the line after $u_{i} \sim N(0, \sigma^{2}z_{i}^{2})$ reads
        ``$u_{i} \sim N(0, \sigma^{2})$''; it is the transformed error and is merged
        with $u_{i}^{*} = u_{i}/z_{i}$. \textbf{Corrected:} the general line is
        written $u_{i} = u/\sqrt{h(z_{i})}$; typeset $u_{i}^{*} = u_{i}/\sqrt{h(z_{i})}$.
        The divided model is written with only the first and last terms
        ($\beta_{1}/\sqrt{h(z_{i})} + \cdots + \beta_{K}X_{Ki}/\sqrt{h(z_{i})}$); the
        $\beta_{2}$ term is \textbf{added}.
  \item Page 3: the heading ``weighted least square'' stands at the top of this page
        but names the procedure that ends page 2, so it heads the last part of
        Section 1. The remark tying it to Lesson 7's $T$ is \textbf{added}. The
        $\Omega$ with $\bar w$ on the diagonal is Lesson 7's $\Omega_{2}$; the name
        ``autocorrelation'' and the section title ``Case 2'' are \textbf{added} (the
        page says only ``Case 1'' for heteroskedasticity). \textbf{Corrected:}
        $u_{t} = \rho u_{t-1} + \varepsilon_{i}$ has subscript $i$; typeset
        $\varepsilon_{t}$. In $u_{t} = \rho Lu_{t} + \varepsilon_{t}$ the page writes
        $u_{t} = \rho Lu_{b}$; read as $u_{t}$.
  \item Page 4: ``Math Recall'' is written without the condition $|a|<1$; it is
        \textbf{added} (the page states $|\rho| < 1$ on the next line).
        \textbf{Corrected:} the moving average is written
        $\varepsilon_{t} + \rho\varepsilon_{t-1} + \rho\varepsilon_{t-2} + \cdots$;
        the third coefficient is $\rho^{2}$. The sentence on $u_{t-1}$ being
        uncorrelated with $\varepsilon_{t}$ is \textbf{added}. \textbf{Corrected:} in
        $u_{t}^{2} = \ldots + 2\rho u_{t}\varepsilon_{t}$ the cross term is
        $2\rho u_{t-1}\varepsilon_{t}$, and the next line has it. The cross term in
        the $\sigma_{u}^{2}$ line is written without the factor $2\rho$; it is
        \textbf{added}; ``the last expectation is zero'' is \textbf{added}.
        \textbf{Corrected:} the following line reads $\sigma_{u}^{2} =
        \rho\sigma_{u}^{2} + \sigma_{\varepsilon}^{2}$; it is $\rho^{2}\sigma_{u}^{2}$,
        as the next line $\sigma_{u}^{2}(1-\rho^{2})$ shows.
  \item Page 5: the top line ``$= \rho\sigma_{u}^{2}$'' finishes $\E(u_{t}u_{t-1})$
        from page 4. \textbf{Corrected:} $u_{t}u_{t-2} = \ldots + \varepsilon u_{t-2}$
        is missing the subscript; typeset $\varepsilon_{t}u_{t-2}$.
        \textbf{Corrected:} the three-lag expectation is written
        $\ldots + \E(\varepsilon_{t}\varepsilon_{t-3})$; it is
        $\E(\varepsilon_{t}u_{t-3})$, which matches the line above it (both are zero).
        The matrix is labelled $\sigma_{u}^{2}\Omega$, $n\times n$, with the note
        ``$|\rho| < 1$\quad $T'T = \Omega^{-1}$''; the note is typeset
        $T'T \propto \Omega^{-1}$, see page 6.
  \item Page 6: the model line is written $u_{t} = \rho u_{t-1}\ \varepsilon_{t}$
        without the plus. ``$/\rho$'' beside equation (2) is read as ``multiply the
        model at $t-1$ by $\rho$''. ``$\varepsilon_{t}$ iid'' with a brace over
        $u_{t} - \rho u_{t-1}$ is kept. The line $y(1-\rho L) = X_{t}'(1-\rho L)\beta
        + \varepsilon_{t}$ is written with $y$ for $y_{t}$. In the column
        $T_{*}y$ the first entry $y_{1} - \rho y_{0}$ is crossed out on the page;
        ``there is no $y_{0}$'' is \textbf{added}. \textbf{Corrected:} the last entry
        is written $y_{n} - y_{n-1}$; it is $y_{n} - \rho y_{n-1}$. The note
        ``$T^{*\prime}T = \Omega^{-1}$'' is replaced by the remark that $T_{*}'T_{*}$ is close to but
        not exactly proportional to $\Omega^{-1}$: dropping the first row changes the
        $(1,1)$ entry.
        The exact $T$ is written with only its first two rows and a diagonal of
        ones; the remaining rows are filled in. The variance check of the first
        observation and $T'T = (1-\rho^{2})\Omega^{-1}$ are \textbf{added} (this is
        the Prais--Winsten transformation; $T_{*}$ alone is Cochrane--Orcutt).
  \item Page 7: \textbf{Corrected:} the lag polynomial is written
        $1 - \rho_{1}L - \rho^{2}L^{2}$; it is $\rho_{2}L^{2}$. \textbf{Corrected:}
        the next line reads $u_{t} - \rho u_{t} - \rho_{2}u_{t-2} = \varepsilon_{t}$;
        it is $\rho_{1}u_{t-1}$. The three lines are kept in the page's order.
  \item Pages of the original, for tracing back: page 1 the divider; page 2 the recap
        of $\Omega$, Case 1 heteroskedasticity, $h(z_{i})$ and the divided model;
        page 3 ``weighted least square'', $\Omega$ with $\bar w$, the AR(1) model and
        the lag operator; page 4 the geometric series, MA($\infty$),
        $\sigma_{u}^{2}$ and the start of $\E(u_{t}u_{t-1})$; page 5 the covariances
        at lags 2 and 3, $\E(u_{t}u_{t-s})$ and $\sigma_{u}^{2}\Omega$; page 6 GLS
        with known $\rho$, $T_{*}$ and the exact $T$; page 7 the AR(2) process.
\end{transcriptnotes}

\end{document}
